\scanpage{39}
\mspage{4}
\uline{NB} Dans cette présentation, symétrie \unsure{s'exprime} par deux aut. involutifs de $H$
\diagram{$\begin{array}{l|ll}
1^\circ)\ g\longmapsto\check g&\varepsilon_0\longmapsto\varepsilon_0^{-1},&\varepsilon_1\longmapsto\varepsilon_1^{-1}\\
(\text{autom. }\uline{\text{ext.}})&\sigma\longmapsto\sigma^{-1}&\\
&\rho_0\longmapsto\rho_1^{-1},&\rho_1\longmapsto\rho_0^{-1}\\
&\rho'_0\longmapsto\rho_1'^{-1},&\rho'_1\longmapsto\rho_0'^{-1}
\end{array}$}
\diagram{$\begin{array}{l|ll}
2^\circ)\ g\longmapsto g^*&\varepsilon_0\longmapsto\varepsilon_1,&\varepsilon_1\longmapsto\varepsilon_0\\
&\sigma\longmapsto\sigma&\\
&\rho_0\longmapsto\rho_1,&\rho_1\longmapsto\rho_0\\
=\operatorname{int}(\sigma)&\rho'_0\longmapsto\rho'_1,&\rho'_1\longmapsto\rho'_0
\end{array}$}
\rule{25mm}{.4pt}

Ici symétrie centrée sur $\sigma$ ou $(\varepsilon_0,\varepsilon_1)$, symétrie binaire $(\sigma^2\in\operatorname{Cent}(H))$.
Si on centre la symétrie sur $\rho\mathrel{\mathop=^{\rm def}}\rho_0$, on trouve symétrie ternaire $(\rho^3\in\operatorname{Cent}(H))$, avec
\[
\varepsilon_0,\ \varepsilon_1=\rho(\varepsilon_0),\quad\varepsilon_\infty=\rho(\varepsilon_1)=\rho^2(\varepsilon_0)
\]
échangés circulairement par $\operatorname{int}(\rho)$,
\diagram{\begin{tikzcd}[column sep=large,row sep=large]
&\varepsilon_\infty\arrow[dl,dash,"\sigma_1=\varepsilon_\infty\varepsilon_0\varepsilon_\infty=\varepsilon_0\varepsilon_\infty\varepsilon_0"']\arrow[dr,dash,"\sigma_0=\varepsilon_1\varepsilon_\infty\varepsilon_1=\varepsilon_\infty\varepsilon_1\varepsilon_\infty"]&\\
\varepsilon_0\arrow[rr,dash,"\sigma_\infty=\varepsilon_0\varepsilon_1\varepsilon_0=\varepsilon_1\varepsilon_0\varepsilon_1"']&&\varepsilon_1
\end{tikzcd}}
\[\left(\begin{array}{c}\Vert\\\sigma\end{array}\right)\]
\[
\rho=(\varepsilon_1\varepsilon_0)^{-1}=(\varepsilon_\infty\varepsilon_1)^{-1}=(\varepsilon_0\varepsilon_\infty)^{-1}
\]
\deleted{$\unsure{\rho=\rho_0\quad(\rho_1)\quad\rho\rho\varepsilon\varepsilon}$}
\margin{\uline{NB} Mais les élémts\par
$\rho_1=(\varepsilon_0\varepsilon_1)^{-1}=\sigma_\infty(\rho)$\par
\deleted{$\unsure{\rho}$}$\ (\varepsilon_1\varepsilon_\infty)^{-1}=\sigma_0(\rho)=\rho(\rho_1)$\par
$(\varepsilon_\infty\varepsilon_0)^{-1}=\sigma_1(\rho)=\rho^2(\rho_1)$\par sont distincts}
\rule{28mm}{.4pt}

Les éléments essentiels du point de vue engendrement de $ST_{11}$ \deleted{sont} semblent être finalement
\[
\varepsilon_0,\varepsilon_1,\sigma,\rho\ (=\rho_0)
\]
plus l'élément
\[\omega=\rho^{-3}=\sigma^2\]
(qui engendre le centre)
\margin{\diagram{$\begin{array}{l}
(\varepsilon_0,\varepsilon_1)\\
(\varepsilon_0,\unsure{\sigma|\rho})\longleftrightarrow(\varepsilon_1,\sigma)\\
(\varepsilon_0,\rho)\longleftrightarrow(\varepsilon_1,\rho)\\
(\sigma,\rho)
\end{array}$}
\unsure{Les quatre expressions de syst. de gén. principaux} de $ST_{11}$.}
\mspage{38}
